\documentclass[a4paper]{article}
\usepackage[pdftex]{graphicx}
\usepackage[utf8]{inputenc}
\usepackage{enumerate}
\usepackage{colortbl}
\usepackage{siunitx}
\sisetup{locale=DE} 
\usepackage{href-ul}
\hypersetup{
	colorlinks=true,
	linkcolor=blue,
	urlcolor=blue}
\usepackage{icomma}
\usepackage{amssymb}
\usepackage{geometry}
\geometry{a4paper, top=15mm, left=15mm, right=15mm, bottom=15mm,
	headsep=10mm, footskip=12mm}

\begin{document}

\begin{enumerate}[1.]

{\bf \item  Aufbau der Atome}

\begin{minipage}{0.65\textwidth}
{\bf Rutherford beschoss eine sehr dünne Goldfolie mit $\alpha$-Teilchen} (Heliumkerne, zweifach positiv geladen) aus einem radioaktiven Präparat und beobachtete, ob und unter welchem Winkel sie abgelenkt werden. Dabei erhielt er die folgenden, überraschenden Ergebnisse:
\end{minipage}
\hfill
\begin{minipage}{0.3\textwidth}
\includegraphics[width=4 cm]{Rutherford129.pdf}
\end{minipage}

\begin{itemize}
\item die meisten $\alpha$-Teilchen gehen ungehindert durch die Folie hindurch.
\item Einige wenige werden unter zum Teil großen Winkeln abgelenkt, sehr wenige werden sogar zurückgestreut.
\end{itemize}

\begin{enumerate}[a)]
	\item Beschreibe den Aufbau eines Atoms und begründe, weshalb der überwiegende Teil der 
	$\alpha$-Teilchen unabgelenkt die Folie durchdringt.
	\item Erkläre das überraschende Phänomen, dass dennoch einige $\alpha$-Teilchen abgelenkt oder gar zurückgestreut werden. Gehe dabei auf die wirkenden Kräfte ein.
\end{enumerate} 

{\bf \item Die konventionelle Schreibweise für den Kern eines Bariumatoms ist ${}^{138}_{~56}{\mathbf{Ba}}$}
\begin{enumerate}[a)]
	\item Gib alle Informationen wieder, die Du aus dieser Schreibweise herauslesen kannst.
	\item Erläutere, weshalb das Bariumatom trotz des positiv geladenen Kerns elektrisch neutral ist.
\end{enumerate}

{\bf \item Halbwertszeit}

\begin{minipage}{0.6\textwidth}
	Eine Gesteinsprobe enthält zum Zeitpunkt $t=0$ $\SI{1,000}{\gram}$ des Thallium-Isotops ${}^{206}_{~81}{\mathbf{Tl}}$ und kein Blei. ${}^{206}_{~81}{\mathbf{Tl}}$  wandelt sich nach einem $\beta$-Zerfall in das stabile Blei um. Die Halbwertszeit hierfür beträgt 4,2 Minuten. 
	\begin{enumerate}[a)]
		\item Stelle die Kernreaktionsgleichung für den $\beta$-Zerfall von Thallium-206 auf.
		\item Fülle die Tabelle aus.
		\item Gibt es einen Zeitpunkt, zu dem alles restlos umgewandelt ist? Begründe!
	\end{enumerate} 
\end{minipage}
\hspace{0.5 cm}
\begin{minipage}{0.4\textwidth}
	\renewcommand{\arraystretch}{2}
	\begin{tabular}{|>{\columncolor[gray]{.8}}p{1.8 cm}|p{1.5 cm}|p{1.5 cm}|}
		\hline
		\rowcolor[gray]{.8}	Zeit in min & Masse Thallium & Masse Blei \\
		\hline
		 0 & $\SI{1000}{\milli\gram}$ & $\SI{0}{\milli\gram}$\\
		\hline
		 4,2 & &\\
		\hline
		 8,4& &\\
		\hline
		 12,6& &\\
		\hline
		 42& &\\
		\hline
	\end{tabular}	
\end{minipage} 

{\bf \item Protonen und Neutronen bestehen aus sogenannten Quarks.}\\ Beschreibe den genauen Aufbau der Nukleonen. Gehe dabei auch auf die elektrische Ladung ein.

{\bf \item Nenne jeweils drei Eigenschaften} von $\alpha$--, $\beta$--, und $\gamma$--Strahlung

{\bf \item Aus einem unbekannten radioaktiven Isotop} ist nach einem $\alpha$--Zerfall das Radon ${}^{218}_{~86}\mathbf{Rn}$ entstanden.\\ Stelle die Kernreaktionsgleichung auf.

{\bf \item Strontium  ${}^{90}_{38}\mathbf{Sr}$ mit einer Halbwertszeit von 29 Jahren} ist ein $\beta$-Strahler und eines der gefähr\-lichsten Folgeprodukte der oberirdischen Kernwaffentests. Nach der fast vollständigen Einstellung der Tests wurde 1966 in Neuherberg bei München eine Sr--90-- Aktivität von $\SI{2,5}{\kilo\becquerel}$ pro Quadratmeter gemessen.

\begin{enumerate}[a)]
	\item Stelle die Kernreaktionsgleichung für den $\beta$-Zerfall von Strontium-90 auf.
	\item Berechne die Aktivität, die pro Quadratmeter im Jahr 2020 zu erwarten ist.
	\item Ermittle das Jahr, ab welchem 75\% zerfallen sind.
\end{enumerate}

%Aufgabe 4
{\bf \item Bei der in der Natur vorkommenden Thorium--Zerfallsreihe} wandelt sich das radioaktive 
${}^{232}_{~90}\mathbf{Th}$ in mehreren Schritten in das stabile Bleiisotop
${}^{208}_{~82}\mathbf{Pb}$ um.\\
Bestimmen Sie durch Rechnung je die Anzahl der dabei auftretenden 
$\alpha$-- und $\beta$--Zerfälle. 

\end{enumerate} 

\fbox{
	\begin{minipage}{0.5\textwidth}
		Zur Lösung bitte \href{https://www.okuyakl.de/phys/p9atoL412/ll412.pdf}{hier klicken} oder den QR-Code scannen.\\
	Weitere Arbeitsblätter gibt es unter 
	
	\href{https://www.okuyakl.de}{www.okuyakl.de}
	\end{minipage}
	\hfill
	\begin{minipage}{0.4\textwidth}
		\includegraphics[width=1.5 cm]{../../viecher/zwe03}
		\includegraphics[width=3 cm]{qr412}
		\includegraphics[width=2 cm]{../../viecher/afanticon1}
		
	\end{minipage}}

\end{document}